\documentclass[border=10pt,crop=true]{standalone}
\usepackage{circuitikz}
% Shared style for 电子学一 Chapter 2 op-amp circuit figures (CircuiTikZ).
\usetikzlibrary{calc}

\ctikzset{
  bipoles/length=0.95cm,
  bipoles/thickness=1,
  resistors/scale=1,
  capacitors/scale=1,
  label distance=2.5mm,
}

\tikzset{
  opfig/.style={font=\small},
  opfigThin/.style={font=\scriptsize},
}

% Geometry (cm) — keep identical across all op-amp schematics
\def\OpInLen{1.55}    % label → input terminal
\def\OpInGap{0.08}    % terminal → wire to op amp +
\def\OpAmpWire{1.00}  % terminal side → op amp + anchor
\def\OpOutLen{1.05}   % op amp out → output terminal
\def\OpJuncL{0.50}    % v_- → feedback junction (left)
\def\OpFbDown{0.55}   % v_- straight down before R_1
\def\OpRvert{1.20}    % R_1 vertical span to ground
\def\OpInSep{0.55}    % dual-input spacing (v_1 / v_2)
\def\OpFbStub{0.40}   % follower: stub right of output before dropping
\def\OpFbClear{0.22}  % follower: below op-amp south for horizontal feedback


\def\OpBusX{-0.90}  % vertical bus: R_3 / R_4 meet here → v_+

\begin{document}
\begin{circuitikz}[american, scale=1.0, transform shape]
  \coordinate (VP) at (\OpBusX, 0);

  \draw (-\OpInLen, \OpInSep) node[left, opfig]{$v_1$} to[short, o-] ++(\OpInGap, 0)
    to[R, l=$R_3$] (\OpBusX, \OpInSep);
  \draw (-\OpInLen, -\OpInSep) node[left, opfig]{$v_2$} to[short, o-] ++(\OpInGap, 0)
    to[R, l=$R_4$] (\OpBusX, -\OpInSep);
  \draw (\OpBusX, \OpInSep) -- (\OpBusX, -\OpInSep);

  \node[opfigThin, anchor=south] at ($(VP)+(0, 0.10)$) {$v_+$};

  \draw (VP) -- ++(\OpAmpWire, 0)
    node[op amp, noinv input up, anchor=+](OA){}
    (OA.-) -- ++(0, -\OpFbDown) coordinate(FB)
    (FB) to[R, l_=$R_1$] ++(0, -\OpRvert) node[ground, scale=0.85]{}
    (FB) to[R, l_=$R_2$] (FB -| OA.out) -- (OA.out)
    to[short, -o] ++(\OpOutLen, 0) node[right, opfig]{$v_o$};
\end{circuitikz}
\end{document}
